\section{Canonical reduction, centering, and endpoint conventions}
\label{ag:sec:normal}

\subsection{The free differential algebra}
Let $\mathscr R=\Q[u_0,u_1,u_2,\ldots]$, where every polynomial uses
finitely many variables. The weight of $u_j$ is $j+1$; its factor degree
is one. The derivation $\D$ increases weight by one and preserves factor
degree. A monomial is \emph{functional} when it is a power of $u_0$,
including $1$, or when its largest positive index occurs at least twice.
Let $\mathscr F$ be the rational span of these monomials.

The following is the normal-form input used in the repository
\cite{Manuscript,HTG}. The terminology and the underlying complement
also occur in the classical differential-algebra construction
\cite[Section~4.4.1]{GRR}. We include a proof to make our coefficient
conventions independent of a particular implementation.

\begin{lemma}[The functional-monomial decomposition]\label{ag:lem:normal}
Every $P\in\mathscr R$ has a unique decomposition
\begin{equation}\label{ag:eq:normal-decomp}
 P=\NF(P)+\D\Prim(P),\qquad \NF(P)\in\mathscr F,
                  \qquad [1]\Prim(P)=0.
\end{equation}
Both maps are rational linear and preserve factor degree. On positive
weight $w$, $\NF(P)$ has weight $w$ and $\Prim(P)$ has weight $w-1$.
\end{lemma}
\begin{proof}
Suppose a monomial has its largest index $r\ge1$ exactly once. Write it
as $u_r u_{r-1}^aV$, where $V$ uses indices below $r-1$. Then
\begin{equation}\label{ag:eq:ibp-step}
 u_r u_{r-1}^aV
 =\D\left(\frac{u_{r-1}^{a+1}V}{a+1}\right)
       -\frac{u_{r-1}^{a+1}\D V}{a+1}.
\end{equation}
Every monomial of the remainder has largest index below $r$. Repeating
this replacement terminates and produces a functional remainder.
Constants of primitives can be set to zero.

For uniqueness, if a nonconstant polynomial $H$ has largest variable
index $j$, the terms $u_{j+1}\partial H/\partial u_j$ in $\D H$ are
nonzero and have largest index $j+1$ exactly once. All other terms have
smaller largest index. They cannot be canceled by a functional
polynomial. Thus $\D H\in\mathscr F$ forces $H$ constant, proving
both directness and uniqueness after the stated normalization.
\end{proof}

This is a vector-space decomposition. For instance,
$\NF(u_1)=0$ but $\NF(u_1^2)=u_1^2$; replacing $\NF(PQ)$ by
$\NF(P)\NF(Q)$ is incorrect. Another useful check is
$\NF(u_0u_4)=u_2^2$. Truncating away $u_4$ before reduction would
lose that term. These are safeguards for the present derivation, not
newly discovered errors in sources which already state the distinction.

\subsection{The centered Gamma representative}
The exponential generating function of \eqref{ag:eq:hierarchy-intro} is
\begin{equation}\label{ag:eq:Bell}
 \A(t)=\sum_{m\ge0}A_m\frac{t^m}{m!}
 =\exp\left(\sum_{j\ge0}\frac{(-1)^j u_jt^{j+1}}{(j+1)!}\right).
\end{equation}
It follows either by the exponential Bell formula or by differentiating
with respect to $t$ and using the defining recurrence. Under the digamma
substitution it is $\Gamma(x)/\Gamma(x-t)$.

Define
\begin{equation}\label{ag:eq:centered}
 \mathcal C(t)=e^{t\D/2}\A(t)
 =\exp\left(\sum_{j\ge0}
       \frac{u_{2j}t^{2j+1}}{2^{2j}(2j+1)!}\right).
\end{equation}
For any formal differential series $P(t)$,
\[
 (e^{t\D/2}-1)P(t)
 =\D\left(\sum_{n\ge1}\frac{t^n\D^{n-1}P(t)}{2^n n!}\right).
\]
Therefore $\NF\mathcal C(t)=\NF\A(t)$. Analytically the centered
representative is $\Gamma(x+t/2)/\Gamma(x-t/2)$, initially near $t=0$
away from Gamma poles.

Translation $u_0\mapsto u_0+c$ commutes with $\D$ and preserves the
functional space. Uniqueness in Lemma~\ref{ag:lem:normal} implies that
it commutes with $\NF$. Thus $\NF\A(t)$ has an Appell factor $e^{tu_0}$.
The centered series also shows factor parity: each exponent term has
odd weight and factor degree one, so every surviving monomial has
factor degree congruent to its weight modulo two.

\begin{lemma}[Appell covariance of the primitive]\label{ag:lem:appell-primitive}
Suppose $P(t)$ and $\NF P(t)$ both satisfy
$\partial_{u_0}P=tP$. If $P-\NF P=\D T$ with $[1]T=0$, then
\[
 \partial_{u_0}T-tT=c(t)
\]
for a scalar series $c(t)$. Consequently, modulo a scalar series,
$T=e^{tu_0}B(t;u_1,u_2,\ldots)$, allowing a harmless Laurent power
of $t$ if necessary.
\end{lemma}
\begin{proof}
The operators $\partial_{u_0}$ and $\D$ commute. Apply
$\partial_{u_0}-t$ to $\D T=P-\NF P$. Its result is zero, so the
kernel statement in Lemma~\ref{ag:lem:normal} makes
$\partial_{u_0}T-tT$ scalar. Solving this elementary equation in $u_0$
gives $T=e^{tu_0}B-c(t)/t$.
\end{proof}
This small observation matters in the affine-symbol proof: after
perturbation, the quadratic part of the primitive has the same
exponential weight as the density and coefficients independent of the
affine coordinate. We do not enlarge the quotient by arbitrary
$x$-dependent primitives.

\subsection{What an analytic substitution does and does not imply}
For every fixed $m$, let $T_m=\Prim(A_m)$. On an interval
$[\alpha,\beta]\subset(0,\infty)$,
\begin{equation}\label{ag:eq:ordinary-gamma-integral}
 \int_\alpha^\beta(-1)^m\Gamma(x)(\Gamma(x)^{-1})^{(m)}\dd x
 =T_m(\psi(\beta),\ldots)-T_m(\psi(\alpha),\ldots)
       +\int_\alpha^\beta N_m(\psi(x),\ldots)\dd x.
\end{equation}
If a fixed-coordinate finite part at zero is used instead, the endpoint
is $T_m(\psi(1),\ldots)-[x^0]T_m(\psi(x),\ldots)$, not zero.
This inherited endpoint convention \cite{Manuscript} is retained.

At positive integral endpoints the entries are generalized harmonic
numbers:
\begin{equation}\label{ag:eq:harmonic-substitution}
 \psi(N+1)=H_N-\gamma,\qquad
 \psi^{(j)}(N+1)=(-1)^{j+1}j!\bigl(\zeta(j+1)-H_N^{(j+1)}\bigr)
 \quad(j\ge1).
\end{equation}
These follow from the standard polygamma recurrence \cite{DLMFpsi}.
Thus the exact normal-form certificates also give finite identities
between Gamma derivatives and harmonic-tail polynomials. They do not
by themselves evaluate every integral of a remaining normal monomial.
